\documentclass{scrartcl}

%This template is based on the DFG rtf form 53_01_de_elan.rtf (see Readme on GitHub for last accessed date)
%
%Author: Martin Hoelzer, hoelzer.martin@gmail.com
%

\usepackage[utf8]{inputenc}
\usepackage[german]{proposal}

% final version?
\setboolean{finalcompile}{false}

% show DFG template version the current LaTeX template is based on in the header. Will be deactivated when 'finalcompile' is set 'true'
\ifthenelse{\boolean{finalcompile}}{}{\ihead*{DFG 53.01 - 03/22}}

% all config is done in Header.tex
\ifthenelse{\boolean{finalcompile}}{
  \usepackage[disable]{todonotes}
  \usepackage[final]{showlabels}
}{
  \usepackage{todonotes}
  %\setlength{\marginparwidth}{2cm}
  \paperwidth=\dimexpr \paperwidth + 6cm\relax
  \oddsidemargin=\dimexpr\oddsidemargin + 3cm\relax
  \evensidemargin=\dimexpr\evensidemargin + 3cm\relax
  \marginparwidth=\dimexpr \marginparwidth + 3cm\relax
  \usepackage{showlabels}
}

\usepackage{lipsum}
\usepackage[pdftex, colorlinks=true, linkcolor=MidnightBlue, citecolor=MidnightBlue, urlcolor=MidnightBlue, pdfauthor={Dr.\ Martin Hoelzer}, pdfpagelayout=SinglePage, bookmarks, bookmarksopen]{hyperref}

\usepackage{gensymb}
\usepackage{graphicx}
\usepackage{wrapfig}

\usepackage{textcomp}
\usepackage{mathtools}
\usepackage{relsize}

\usepackage[numbered]{bookmark}

%\usepackage{titlesec}
%\usepackage{placeins}

\usepackage[font=small,labelfont={bf,it},textfont=it]{caption}
\captionsetup{format=plain}

\usepackage[noabbrev]{cleveref}
% see also: https://www.namsu.de/Extra/pakete/Cleveref.html

% make numbers look better in funding section
%siunitx package is loaded already in proposal.sty
\sisetup{detect-all=true}

\graphicspath{{figures/}}
\addbibresource{literature.bib}
% use \addbibresource again to add multiple bib files
\AtBeginBibliography{\small}

% command \todo now deprecated by using todonotes
%\newcommand{\todo}[1]{\xspace{\textcolor{red}{\bfseries[TODO: #1]}}\xspace}
\newcommand{\cp}[1]{\xspace{\textcolor{blue}{\bfseries[COPIED: #1]}}\xspace}

\usepackage{tikz}
\usepackage{lipsum}
\usepackage{enumitem}
\xdef\Colored{0}
\ifnum\Colored > 0
\xdef\LWS{0.6mm}
\xdef\ColorYO{yellow!30!orange}
\xdef\ColorYR{yellow!70!red}
\xdef\ColorRY{red!90!yellow}
\xdef\ColorYB{yellow!90!brown}
\xdef\ColorYRB{yellow!90!red}
\xdef\ColorR{red}
\xdef\ColorO{orange}
\else
\xdef\LWS{0.4mm}
\xdef\ColorYO{gray!80!black}
\xdef\ColorYR{gray!70}
\xdef\ColorRY{gray!60!black}
\xdef\ColorYB{gray!90}
\xdef\ColorYRB{gray!50}
\xdef\ColorR{black}
\xdef\ColorO{gray!95!black}
\fi
\newenvironment{enuma}{\begin{enumerate}[label={\alph{enumi})}, itemsep=-1mm]}{\end{enumerate}}
\newlength\tipWidth
\newlength\tipSeparator
\newlength\tipTextWidth
\newlength\addedWidth
\newlength\addedHeight
\newlength\RectHeight
\newlength\testL
\newlength\lengthA
\newlength\lengthB
\newlength\lengthC
\setlength\addedHeight{3mm}
\setlength\addedWidth{4mm}
\setlength\tipSeparator{6mm}
\newsavebox{\lambBox}
\newsavebox{\lambVBox}
\newsavebox{\tipBox}
\newsavebox{\tipVBox}
\newcommand{\myTip}[4][\tipWidth]{%
\def\alig{#2}%
\def\aligm{m}%
\def\aligt{t}%
\savebox{\lambBox}{\hbox{\begin{tikzpicture}[inner sep=0,outer sep=0]
\fill[rounded corners=0.35mm,rotate=-15,gray!90!black] (0.05,-0.15)rectangle(-0.1,-0.02);
\draw[line width=0.9mm] (0.2,0.01) to[in=330 , out=170] (-0.2,0.11);
\draw[line width=1mm] (0.3,0.15) to[in=330 , out=170] (-0.25,0.25);
\draw[line width=1.1mm] (0.35,0.3) to[in=350 , out=180] (-0.27,0.36);
\draw[line width=\LWS, color=\ColorYO] (-0.15,0.42) to[in=310,out=80] (-0.3,0.8)to[in=215,out=130](-0.37,1.15) to[out=45,in=150] (0.2,1.35)to[in=110,out=-30](0.56,0.99)to[in=45,out=290](0.41,0.69)to[in=80,out=215](0.2,0.37);
\draw[line width=0.2mm] (0.,0.33)--(-0.02,0.72);
\fill[\ColorYR] (-0.02,0.72)--(0.0,0.93)--(0.055,0.79)--(0.07,0.9)--(0.13,0.79)--(0.28,0.8)--(0.16,0.61)--(0.09,0.65)--(0.07,0.5)--(0.046,0.54)--cycle;
\fill[\ColorRY] (-0.02,0.72)--(0.008,0.78)--(0.05,0.68)--(0.07,0.84)--(0.15,0.74)--(0.28,0.8)--(0.165,0.69)--(0.10,0.74)--(0.07,0.63)--(0.045,0.63)--(0.015,0.72)--cycle;
\draw[line width=0.2mm] (0.02,0.3)--(0.28,0.8);
\draw[line width=0.7mm,\ColorYB](0.25,1.46)--(0.32,1.69) (0.41,1.33)--(.61,1.57) (0.54,1.19)--(.7,1.33) (0.03,1.5)--(-0.0,1.8) (-0.18,1.42)--(-0.3,1.62) (-0.4,1.3)--(-0.58,1.49);
\draw[line width=0.3mm,\ColorYRB] (-0.3,1.05)to[in=120,out=20](0.5,0.96);
\end{tikzpicture}}}%
\savebox\lambVBox{\vbox{\usebox{\lambBox}}}%
\setlength\tipTextWidth{0 pt}%
\addtolength\tipTextWidth{#1}%
\addtolength\tipTextWidth{-\wd\lambBox}%
\addtolength\tipTextWidth{-\tipSeparator}%
\addtolength\tipTextWidth{-\addedWidth}%
\addtolength\tipTextWidth{-\addedWidth}%
\savebox{\tipBox}{\hbox{\begin{minipage}[inner sep=0,outer sep=0]{\tipTextWidth}{\noindent\bfseries WP#3 $\vert$} #4\end{minipage}}}%
\savebox{\tipVBox}{\vbox{\usebox{\tipBox}}}%
\setlength\testL{\ht\tipVBox}%
\addtolength\testL{\dp\tipVBox}%
\addtolength\testL{-\ht\lambBox}%
\addtolength\testL{-\dp\lambBox}%
\setlength\lengthA{\ht\tipVBox}%
\addtolength\lengthA{-0.5\ht\lambBox}%
\setlength\lengthB{0.5\ht\lambBox}%
\addtolength\lengthB{-\ht\tipVBox}%
\setlength\lengthC{0.5\ht\tipVBox}%
\addtolength\lengthC{0.5\dp\tipVBox}%
\begin{tikzpicture}[inner sep=0, outer sep =0]%
\node[inner sep=0,outer sep=0] at ({-0.5\wd\lambBox-0.5\tipSeparator},{\ifx\alig\aligm 0\else \ifdim\testL > 0pt \lengthA\else 0\fi\fi}) {\usebox{\lambBox}};
\node[inner sep=0,outer sep =0] at ({0.5\tipTextWidth+0.5\tipSeparator},{\ifx\alig\aligm 0\else \ifdim\testL > 0pt 0\else \lengthB\fi\fi}) {\usebox{\tipBox}};
\draw[\ColorO,line width=0.5mm] (0,\ifdim\testL> 0pt -\ht\tipVBox \else -0.5\ht\lambBox\fi)--(0,\ifdim\testL> 0pt \lengthC \else 0.5\ht\lambBox\fi);
\draw[\ColorR,line width=0.5mm,rounded corners=1mm] ({-\wd\lambBox-\addedWidth-\tipSeparator/2},{\ifdim\testL >0pt -\ht\tipVBox\else -0.5\ht\lambBox\fi-\addedHeight})rectangle({\addedWidth+\tipSeparator/2+\tipTextWidth},{\ifdim\testL > 0pt  \lengthC\else 0.5\ht\lambBox\fi+\addedHeight});
\end{tikzpicture}%
}



%%%%%%%%%%%%%%%%%%%%%%%%%%%%%%%%%%%%%%%%%%%%%%%%%%%%%%%%%%%%%%%%%%%%%%%%%%%%%
%%%%  TITELSEITE  %%%%%%%%%%%%%%%%%%%%%%%%%%%%%%%%%%%%%%%%%%%%%%%%%%%%%%%%%%%
%%%%%%%%%%%%%%%%%%%%%%%%%%%%%%%%%%%%%%%%%%%%%%%%%%%%%%%%%%%%%%%%%%%%%%%%%%%%%

\newcommand{\applicants}{Dr.\ Martin H\"olzer, Robert Koch Institute, Berlin}
\newcommand{\project}{Project title}

% declare bibliography categroies

\addtocategory{reviewed}{Hoelzer:16} 
\nocite{}
\addtocategory{nonreviewed}{Desiro:18} 
\addtocategory{patents_pending}{} 
\addtocategory{patents}{} 
% 
\begin{document}
\pagestyle{empty}
\setcounter{page}{0}

{\selectlanguage{german}

  \vspace*{20ex}

  \begin{center}
  \large

  {\bfseries Antrag an die Deutsche Forschungsgemeinschaft}

  Kennedyallee 40, 53175 Bonn

  \vspace{4ex}

  zur Förderung des Forschungsprojekts

  \vspace{9ex}

  \begin{bfseries}
    This is a fancy project title to rule the world
  \end{bfseries}

  \medskip

	\todo{Erstantrag} Projektantrag

  \vspace{9ex}

	Dr.\ Martin H\"olzer

 	Robert Koch Institute
  
 	Method Development and Research Infrastructure

  Bioinformatics (MF1)

	Nordufer 20, 13353 Berlin
	 
	\vspace{9ex}

	\today
  \end{center}
}

%%% Local Variables:
%%% mode: latex
%%% TeX-master: "dfg-german"
%%% End:


%%%%%%%%%%%%%%%%%%%%%%%%%%%%%%%%%%%%%%%%%%%%%%%%%%%%%%%%%%%%%%%%%%%%%%%%%%%%%
%%%%  MAIN MATTER  %%%%%%%%%%%%%%%%%%%%%%%%%%%%%%%%%%%%%%%%%%%%%%%%%%%%%%%%%%
%%%%%%%%%%%%%%%%%%%%%%%%%%%%%%%%%%%%%%%%%%%%%%%%%%%%%%%%%%%%%%%%%%%%%%%%%%%%%
\cleardoublepage
\pagestyle{plain}

%%%%%%%%%%%%%%%%%%%%%%%%%%%%%%%%%%%%%%%%%%%%%%%%%%%%%%%%%%%%%%%%%%%%%%%%%%%%%
%%%% For first-time proposals, please also note the following: Please provide a
% concise description of your previous postgraduate training and research,
% highlighting your abilities and demonstrating that you are capable of carrying
% out the proposed project. Describe your scientific track record to date,
% mentioning research topics that you’ve worked on so far. Previous scientific
% accomplishments must not be directly related to the project.

%%%%%%%%%%%%%%%%%%%%%%%%%%%%%%%%%%%%%%%%%%%%%%%%%%%%%%%%%%%%%%%%%%%%%%%%%%%%%%%
%%%%  PROJECT DESCRIPTION - PROJECT PROPOSALS  %%%%%%%%%%%%%%%%%%%%%%%%%%%%%%%%
%%%%%%%%%%%%%%%%%%%%%%%%%%%%%%%%%%%%%%%%%%%%%%%%%%%%%%%%%%%%%%%%%%%%%%%%%%%%%%%
%% Sections 1-4 must not exceed 17 pages in total.
{\raggedright{} \normalsize \bfseries 
	Beschreibung des Vorhabens - Projektanträge \par
	\applicants{[Name, Vorname, Ort aller Antragstellenden]} \par
	\project{[Titel des Projekts]} \par
	\rule{\textwidth}{0.5pt} \par
	%Project Description
}

%%%%%%%%%%%%%%%%%%%%%%%%%%%%%%%%%%%%%%%%%%%%%%%%%%%%%%%%%%%%%%%%%%%%%%%%%%%%%
%%%%  STATE OF THE ART AND PRELIMINARY WORK %%%%%%%%%%%%%%%%%%%%%%%%%%%%%%%%%
%%%%%%%%%%%%%%%%%%%%%%%%%%%%%%%%%%%%%%%%%%%%%%%%%%%%%%%%%%%%%%%%%%%%%%%%%%%%%
\section{Ausgangslage}
\label{sec:work-report}

\subsection{Stand der Forschung und eigene Vorarbeiten}
% For renewal proposals, please report on your previous work. This report
% should also be understandable without referring to additional literature.

This is a dummy citation~\cite{Hoelzer:17}. Yeah. And another
one~\cite{Gerst:18}. Wuhuh. These citations~\cite{Hoelzer:16, Desiro:18} will
not be shown in the "\ref{sec:bib}~Bibliography" section, because they are
already listed as \verb=\addtocategory{reviewed}{Hoelzer:16}= and
\verb=\addtocategory{nonreviewed}{Desiro:18}= in the praeambl and therefore
shown in subsubsections \emph{Articles published by outlets with scientific
quality \dots} and \emph{Other publications, both peer-reviewed and non-peer-reviewed}. 

\lipsum[1]


\subsection{Projektbezogenes Publikationsverzeichnis Ihrer Arbeiten}
%Please include a list of own publications that are related to the proposed project. 
%Sections 1.2.1 and 1.2.2 together must not exceed 10 publications; please number them consecutively.

%a maximum of ten publications
\subsubsection{Veröffentlichte Arbeiten aus Publikationsorganen mit
  wissenschaftlicher Qualitätssicherung, Buchveröffentlichungen sowie
  bereits zur Veröffentlichung angenommene, aber noch nicht
  veröffentlichte Arbeiten}

% Own literature can be designated with a prefix, default 'E' for "Eigene". 
% If the function is not desired, delete the "labelprefix" commands directly below this comment 
% AND in the section: "Literaturverzeichnis zum Stand der Forschung, zu den Zielen und dem Arbeitsprogramm"
% Then, all your references should be numbered consecutively.
\newrefcontext[labelprefix=E]
\printbibliography[category=reviewed, heading=none, env=bibliographyNUM, resetnumbers]

\subsubsection{Andere Veröffentlichungen mit und ohne wissenschaftliche Qualitätssicherung}
\printbibliography[category=nonreviewed, heading=none, env=bibliographyNUM, resetnumbers=false]

\subsubsection{Patente}

\subsubsubsection{Angemeldet}
\printbibliography[category=patents_pending, heading=none]
None.

\subsubsubsection{Erteilt}
\printbibliography[category=patents, heading=none]
None.

%%%%%%%%%%%%%%%%%%%%%%%%%%%%%%%%%%%%%%%%%%%%%%%%%%%%%%%%%%%%%%%%%%%%%%%%%%%%%
%%%%  OBJECTIVES AND WP %%%%%%%% %%%%%%%%%%%%%%%%%%%%%%%%%%%%%%%%%%%%%%%%%%%%
%%%%%%%%%%%%%%%%%%%%%%%%%%%%%%%%%%%%%%%%%%%%%%%%%%%%%%%%%%%%%%%%%%%%%%%%%%%%%
\section{Ziele und Arbeitsprogramm}

\subsection{Voraussichtliche Gesamtdauer des Projekts}
[Text]

\subsection{Ziele}
\let\oldpara=\theparagraph
\addtocounter{secnumdepth}{1}
\renewcommand{\theparagraph}{Ziel \arabic{paragraph}}

\paragraph{\textnormal{Goal one is ...}}

\subsection{Arbeitsprogramm inkl. vorgesehener Untersuchungsmethoden}
%For each applicant
\todo[inline]{ca.\ 6-8 pages; ca. 4-8 WPs}

%%%%%%%%%%%%%%%%%%%%%%%%%%%%%%%%%%%%%%%%%%%%%%%%%%%%%%%%%%%%%%%%%%%%%%%% 
%%%%% BEGIN OF WORK PACKAGES %%%%%%%%%%%%%%%%%%%%%%%%%%%%%%%%%%%%%%%%%%%
%%%%%%%%%%%%%%%%%%%%%%%%%%%%%%%%%%%%%%%%%%%%%%%%%%%%%%%%%%%%%%%%%%%%%%%% 
%\titlespacing*{\paragraph}{0pt}{-0.2\baselineskip}{0.2\baselineskip}
\addtocounter{secnumdepth}{1}
%\newcommand\WP{WP\arabic{paragraph}}
\renewcommand{\theparagraph}{AP\arabic{paragraph}}
\RedeclareSectionCommand[beforeskip=-2.5ex plus-1ex minus -.2ex,%
afterskip=2.5ex plus 0.2ex minus 0.2ex]{paragraph}
%numbered sub-workpackages are supported - just use \subparagraph{}
\RedeclareSectionCommand[indent=0pt, beforeskip=-3.5ex plus-1ex minus -.2ex,%
afterskip=2.5ex plus 0.2ex minus 0.2ex]{subparagraph}
%%%%%%%%%%%%%%%%%%%%%%%%%%%%%%%%%%%%%%%%%%%%%%%%%%%%%%%%%%%%%%%%%%%%%%%%%%%%%

Für jede Antragstellerin und jeden Antragsteller

\hrulefill
\paragraph{Erstes Arbeitspaket.}
\label{wp:1}

\lipsum[1]

\subsubsection*{Investigation of unique characteristics of the \emph{Myotis} immune system}
\lipsum[2]

\bigskip
% I am sorry, you have to change the wp number here manually {1}
% but hey, you get a nice summary box!
\noindent\myTip[1.0\textwidth]{t}{1}{\emph{\vspace*{0.1cm}\textbf{Zusammenfassung}\newline
\textbf{Ziel:}
In this work package,  \newline
\textbf{Methode:}
Based on our unique  \newline
\textbf{Ergebnis:}
\lipsum[1]
}}



\hrulefill
\paragraph{Zweites Arbeitspaket.}
\label{wp:2}

\lipsum[1]

\subsubsection*{Investigation of unique characteristics of the \emph{Myotis} immune system}
\lipsum[2]

\bigskip
% I am sorry, you have to change the wp number here manually {1}
% but hey, you get a nice summary box!
\noindent\myTip[1.0\textwidth]{t}{2}{\emph{\vspace*{0.1cm}\textbf{Zusammenfassung}\newline
\textbf{Ziel:}
In this work package,  \newline
\textbf{Methode:}
Based on our unique  \newline
\textbf{Ergebnis:}
\lipsum[1]
}}


%%%%%%%%%%%%%%%%%%%%%%%%%%%%%%%%%%%%%%%%%%%%%%%%%%%%%%%%%%%%%%%%%%%%%%%%%%%%%
%%%%  END OF WORK PACKAGES  %%%%%%%%%%%%%%%%%%%%%%%%%%%%%%%%%%%%%%%%%%%%%%%%%
%%%%%%%%%%%%%%%%%%%%%%%%%%%%%%%%%%%%%%%%%%%%%%%%%%%%%%%%%%%%%%%%%%%%%%%%%%%%%


%%%%%%%%%%%%%%%%%%%%%%%%%%%%%%%%%%%%%%%%%%%%%%%%%%%%%%%%%%%%%%%%%%%%%%%%%%%%%
%%%%  TIMELINE  %%%%%%%%%%%%%%%%%%%%%%%%%%%%%%%%%%%%%%%%%%%%%%%%%%%%%%%%%%%%%
%%%%%%%%%%%%%%%%%%%%%%%%%%%%%%%%%%%%%%%%%%%%%%%%%%%%%%%%%%%%%%%%%%%%%%%%%%%%%
%% This is not directly part of the DFG TRF template, but always helpful in my eyes
\let\theparagraph=\oldpara
\paragraph*{Zeitplan}
\vspace{-0.5cm}
\begin{figure}[h]
	\centering
  	\resizebox{\textwidth}{!}{\includestandalone[mode=tex]{gantt/gantt}}
	\caption{Arbeitspakete Zeitplan in Monaten (M) und zugeordneten Personenmonaten (PM) f\"ur Institute of ABC (ABC) und Institute for CDE (CDE)}
	\label{fig:timeline}
\end{figure}

\subsection{Umgang mit Forschungsdaten}
Data generated during this project will be used for scientific publications in
preprint and peer-reviewed journals. Therefore, all necessary data (such as raw
sequencing data) will be deposited in publicly available repositories
(\href{https://www.ncbi.nlm.nih.gov/geo/}{NCBI GEO} and
\href{https://www.ncbi.nlm.nih.gov/sra}{SRA}, \href{https://osf.io/}{OSF}). To
allow complete reproducibility of all analyses, the source code and executed
comands will be also made publicly available
(\href{https://github.com/}{GitHub}).  Thus, all relevant data will become
accessible for future use. 
\todo[inline]{FAIR principals}

\subsection{Relevanz von Geschlecht und/oder Vielfältigkeit}
\todo[inline]{Text}

\section{Literaturverzeichnis zum Stand der Forschung, zu den Zielen und dem Arbeitsprogramm}
\label{sec:bib}
%\AtNextBibliography{\small}
\newrefcontext[labelprefix= ]
\printbibliography[notcategory=reviewed, notcategory=nonreviewed, notcategory=patents_pending, notcategory=patents, heading=none, env=bibliographyNUM]

\backmatter
%%%%%%%%%%%%%%%%%%%%%%%%%%%%%%%%%%%%%%%%%%%%%%%%%%%%%%%%%%%%%%%%%%%%%%%%%%%%%%%
%%%%  SUPPLEMENT  (new since 04/2020) %%%%%%%%%%%%%%%%%%%%%%%%%%%%%%%%%%%%%%%%%
%%%%%%%%%%%%%%%%%%%%%%%%%%%%%%%%%%%%%%%%%%%%%%%%%%%%%%%%%%%%%%%%%%%%%%%%%%%%%%%
\section{Begleitinformationen zum Forschungskontext}
% Section 5 et seq. must not exceed 10 pages.

\subsection{Angaben zu ethischen und/oder rechtlichen Aspekten des Vorhabens}

\subsubsection{Allgemeine ethische Aspekte}
\todo[inline]{Text}

\subsubsection{Erläuterungen zu den vorgesehenen Untersuchungen am Menschen, an vom Menschen entnommenem Material oder mit identifizierbaren Daten}
\todo[inline]{Text}

\subsubsection{Erläuterungen zu den vorgesehenen Untersuchungen bei Versuchen an Tieren}
\todo[inline]{Text}

\subsubsection{Erläuterungen zu Forschungsvorhaben an genetischen Ressourcen (oder 
darauf bezogenem traditionellem Wissen) aus dem Ausland}
\todo[inline]{Text}

\subsubsection{Erläuterungen zu möglichen sicherheitsrelevanten Aspekten („Dual-Use 
Research of Concern``, Außenwirtschaftsrecht)}
\todo[inline]{Text}

%%%%%%%%%%%%%%%%%%%%%%%%%%%%%%%%%%%%%%%%%%%%%%%%%%%%%%%%%%%%%%%%%%%%%%%%%%%%%%%
%%%%  PEOPLE/COLLAB/FUNDING  %%%%%%%%%%%%%%%%%%%%%%%%%%%%%%%%%%%%%%%%%%%%%%%%%%%
%%%%%%%%%%%%%%%%%%%%%%%%%%%%%%%%%%%%%%%%%%%%%%%%%%%%%%%%%%%%%%%%%%%%%%%%%%%%%%%

\subsection{Angaben zur Dienststellung}
%for each applicant, state the last name, first name, and employment status
%(including duration of contract and funding body, if on a fixed-term contract)
H\"olzer, Martin -- postdoctoral researcher, \todo[inline]{additional information}

\subsection{Angaben zur Erstantragstellung}
%only if applicable: Last name, first name of first-time applicant
H\"olzer, Martin.

\subsection{Zusammensetzung der Projektarbeitsgruppe}
%List only those individuals who will work on the project but will not be paid
%out of the project funds. State each person's name, academic title, employment
%status, and type of funding. 
The following individuals will work on the project but will not be paid out of
the project funds:

\begin{itemize}
\item Dude1
\item Dude2
\end{itemize}

\subsection{Zusammenarbeit mit Wissenschaftlerinnen und Wissenschaftlern in Deutschland in diesem Projekt}
\todo[inline]{Text}

\subsection{Zusammenarbeit mit Wissenschaftlerinnen und Wissenschaftlern im Ausland in diesem Projekt}
\todo[inline]{Text}

\subsection{Wissenschaftlerinnen und Wissenschaftler, mit denen in den letzten 
drei Jahren wissenschaftlich zusammengearbeitet wurde}
% This information will help avoid potential conflicts of interest.
\begin{itemize}
  \item Prof.\ Dr.\ Dr.\ Foo Bar
\end{itemize}

\subsection{Projektrelevante Zusammenarbeit mit erwerbswirtschaftlichen Unternehmen}
%If applicable, please note the EU guidelines on state aid or contract your research institution in this regard.
None.

\subsection{Projektrelevante Beteiligungen an erwerbswirtschaftlichen Unternehmen}
%Information on connections between the project and the production branch of the enterprise
None.

\subsection{Apparative Ausstattung}
%List larger instruments that will be available to you for the project. These may include large computer facilities if computing capacity will be needed.
\todo[inline]{Text}
Man kann etwas in Abbildung~\ref{fig:some_nice_graph} sehen\todo[inline]{the figure has to be made}.
\begin{figure}
\centering
\missingfigure{Fancy Laborausstattung}
\caption{Tolle Laborausstattung}
\label{fig:some_nice_graph}
\end{figure}

\subsection{Weitere Antragstellungen}
%List any funding proposals for this project and/or major instrumentation previously submitted to a third party.
\todo[inline]{Text}

\subsection{Weitere Angaben}
%Please use this section for any additional information you feel is relevant which has not been provided elsewhere.
In submitting a proposal to the DFG, I agree to adhere to the DFG's rules of good scientific practice and the \href{https://www.nature.com/articles/sdata201618}{FAIR principles}.

%\section{Requested modules/funds}
%Explain each item for each applicant (stating last name, first name)

%%%%%%%%%%%%%%%%%%%%%%%%%%%%%%%%%%%%%%%%%%%%%%%%%%%%%%%%%%%%%%%%%%%%%%%%%%%%%%%
%%%%  REQUESTED MODULES/FUNDS  %%%%%%%%%%%%%%%%%%%%%%%%%%%%%%%%%%%%%%%%%%%%%%%%%%%
%%%%%%%%%%%%%%%%%%%%%%%%%%%%%%%%%%%%%%%%%%%%%%%%%%%%%%%%%%%%%%%%%%%%%%%%%%%%%%%
\section{Beantragte Module/Mittel}
% Explain each item for each applicant (stating last name, first name).

\subsection{Basismodul}

\subsubsection{Personalmittel}
\begin{funds}[funding for staff]
The following staff positions are requested for \todo[inline]{xx} years each:

\positionmul{Research associate, TV-L 13, 36 months}{5375}{36}
\positionmul{Hiwi, TV-L 13, 12 months}{450}{12}

\end{funds}


\subsubsection{Sachmittel}
\begin{funds}[direct project costs]

\subsubsubsection{Geräte bis 10.000 Euro, Software und Verbrauchsmaterial}

\position{Next-Generation Sequencing, Illumina HiSeq\,2500}{3000}

\subsubsubsection{Reisemittel}
We apply for a total of 10\,000\,\euro\ for travel expenses.
\position{Travel}{10000}

\subsubsubsection{Mittel für wissenschaftliche Gäste \textnormal{(ausgenommen Mercator-Fellow)}}
None.

\subsubsubsection{Mittel für Versuchstiere}
None.

\subsubsubsection{Sonstige Mittel}
None.

\subsubsubsection{Publikationsmittel}
We apply for a total of 1500\euro\ for publication expenses (750\,\euro\ per
year). Furthermore, we will submit articles to open access preprint repositories
such as \href{https://www.biorxiv.org/}{\textit{bioRxiv}}.
\position{Publication costs}{1500}

\end{funds}

\subsubsection{Investitionsmittel}

\subsubsubsection{Geräte über 10.000 Euro}
Our laboratory is well equipped; all necessary instruments are available.

\subsubsubsection{Großgeräte über 50.000 Euro}
Our laboratory is well equipped; all necessary instruments are available.


%%%%%%%%%%%%%%%%%%%%%%%%%%%%%%%%%%%%%%%%%%%%
%% The following are additional/other modules one might want to apply for
%% otherwise, just delete/comment
%%%%%%%%%%%%%%%%%%%%%%%%%%%%%%%%%%%%%%%%%%%%

\vspace*{2cm}\todo[inline]{The following are additional/other modules one might want to apply for. Otherwise, just delete/comment}

\subsection{Modul Eigene Stelle}
\todo[inline]{Text}

\subsection{Modul Vertretung}
\todo[inline]{Text}

\subsection{Modul Rotationsstellen}
\todo[inline]{Text}

\subsection{Modul Mercator Fellow}
\todo[inline]{Text}

\subsection{Modul Projektspezifische Workshops}
\todo[inline]{Text}

\subsection{Modul Öffentlichkeitsarbeit}
\todo[inline]{Text}

\subsection{Modul Pauschale für Chancengleichheitsmaßnahmen}
\todo[inline]{Bitte erläutern Sie ausführlich die geplanten Maßnahmen zur Förderung der Vielfalt und Chancengleichheit. Sofern Sie Ihren Sachbeihilfeantrag innerhalb eines Verbundkontextes stellen, so kann diese Pauschale nur gebündelt im Koordinierungsprojekt beantragt werden.}

\section{Unterschrift}

\includegraphics[width=0.18\textwidth]{/signatures/signature.pdf}\\
Dr.\ Martin H\"olzer

\section{Anhänge}
\begin{itemize}
  \item Curriculum Vit\ae, Dr.\ Martin H\"olzer
\end{itemize}

\end{document}

%%% Local Variables:
%%% mode: latex
%%% TeX-master: t
%%% End:
